\chapter{Results}
\label{ch:results}

% RULES FOR THIS CHAPTER:
% 1. NEVER hand-copy a number. Tables are \input from thesis/tables/ and figures
%    \includegraphics from thesis/figures/, both generated from the repository
%    (MLflow runs + eval reports). See thesis/README.md.
% 2. Every number must be labelled as either DIAGNOSTIC (in-sample; used while
%    building the apparatus) or SEALED (held-out split; reportable). An in-sample
%    number may never be presented as the system's accuracy.
% 3. Sections marked [PENDING] wait on the reader-selection run; they carry a
%    visible placeholder so a missing result can never be mistaken for a null one.

\section{Overview of the Result Set}
\label{sec:res-overview}
% TODO: one page orienting the reader: which experiments exist, which question each
% answers, and which are diagnostic vs sealed.

\section{Model Cohort}
\label{sec:res-cohort}
% TODO: the evaluated open-weights local document \ac{VLM} cohort and the selection
% rationale (architecture families, size classes, licence, German capability).

\section{Reading and Structuring: Cohort Comparison}
\label{sec:res-cohort-comparison}
% TODO: per-model field \ac{F1} over the corpus. Table: thesis/tables/cohort.tex

\section{Two-Stage versus Single-Shot Architecture}
\label{sec:res-architecture}
% TODO: reader+structurer against one model going image -> structured output;
% per-field failure profiles. This answers sub-question 2 of Sec.~\ref{sec:research-questions}.

\section{Structured-Output Prompting}
\label{sec:res-structured}
% TODO: free-form transcript plus deterministic adapter vs. asking the model for
% \ac{JSON} directly; include the honesty axis (does the model invent values for
% fields that are absent from the page?).

\section{Efficiency on Local Hardware}
\label{sec:res-efficiency}
% TODO: decode throughput and peak memory per model; which models fit the local
% memory envelope and which do not; the accuracy-versus-cost frontier.
% Answers sub-question 3.

\section{Measurement Validity}
\label{sec:res-measurement}
% TODO: the chain of ground-truth and scoring corrections, and how much apparent
% model error each one turned out to be. This is a result in its own right: a
% sizeable part of the initially observed error was scoring artefact, not model
% failure. Present it as a table of correction -> effect.

\section{Error Attribution: Reading versus Structuring}
\label{sec:res-attribution}
% TODO: the decisive experiment. Same structuring model, two inputs: real
% transcripts vs reference-quality transcripts. The gap between the two bounds how
% much of the end-to-end error is caused by reading rather than structuring.
% Report per-field-group decomposition and the share of missing values that are
% literally absent from the transcript.
% Table: thesis/tables/attribution.tex
% Answers sub-question 4 -- and motivates the reader-selection work below.

\section{Reader Selection}
\label{sec:res-reader-selection}
% [PENDING] Waits on the reader comparison run (rented GPU session).
\noindent\textit{[Pending --- reader-comparison results are not yet available.
This section will report the candidate comparison and the selected reader. No
number is stated here until the run completes.]}

\section{Re-Baseline with the Selected Reader}
\label{sec:res-rebaseline}
% [PENDING] Waits on transcript regeneration with the selected reader.
\noindent\textit{[Pending --- the end-to-end score after replacing the reader will
be reported here.]}

\section{Structurer Adaptation}
\label{sec:res-finetune}
% [PENDING, CONDITIONAL] Low-rank adaptation of the structurer runs only if the
% re-baseline stays below the target. If it is not needed, that is itself a
% reportable result -- say so explicitly rather than omitting the section.
\noindent\textit{[Pending --- conditional on the re-baseline. If adaptation proves
unnecessary, this section reports that outcome.]}
